\begin{abstract}
CDNs gate protected objects with signed URLs and cookies, using either an HMAC key shared with every edge or RSA, ECDSA or Ed25519, which a quantum computer would break. Replacing them with per-token ML-DSA or SLH-DSA signatures yields tokens of \num{3350} to \num{22907} bytes that exceed URL budgets and inflate every segment request. We present EMBER, a CDN profile of Merkle batch signatures: the issuer places all tokens of an epoch in a salted, domain-separated Merkle tree, signs only the root, and each token carries an authentication path. We contribute a threat model stating that HMAC is already post-quantum and that the gain is removing minting secrets from edges; root-dissemination modes for an edge fleet, including root stapling; a security reduction sketch; a measurement study against 13 baselines, among them multi-MAC, delegated Ed25519, KEM-wrapped MAC keys, MTL mode and cached ML-DSA; and a fleet simulator validated against a Go prototype. With $2^{16}$ tokens per epoch an EMBER token is 503~B, against \num{3350}~B for ML-DSA-44. Warm verification takes 2.34\us{} natively, 43$\times$ less than ML-DSA-44 verification but 7.9$\times$ HMAC, and one issuer core produces 1.28 million tokens/s. Over HTTP/2 a session needs 0.144$\times$ the header bytes of split-cookie ML-DSA-44. Stapling lowers the share of simulated requests that wait for a root to 0.049\%. We also report negative results. An adaptive epoch controller performed no better than a fixed 10~ms epoch. Period roll-up removed only 1.4--2.2\% of root misses. A hybrid Ed25519+ML-DSA root raised simulated edge CPU per request by 41.6\%.
\end{abstract}

\begin{IEEEkeywords}
Access tokens, content delivery networks, Merkle trees, ML-DSA, post-quantum cryptography, signed URLs, SLH-DSA.
\end{IEEEkeywords}
